\subsection{Splitting transfers and rigidified Picard norms}\label{sec:paper-splitting}
In this subsection, we introduce the splitting transfer subgroup $I_F(D)$ and the quotient $Q_F(D)=K_2(D)/I_F(D)$, which measures how far $K_2(D)$ is from being generated by transfers from splitting fields.
The main result is Proposition~\ref{prop:paper-picard-eval}: over a field complete for a nontrivial rank-$1$ nonarchimedean valuation, the evaluations at rational points of a $K_2$-class on a smooth affine scheme are locally constant modulo the transfers from a fixed finite separable splitting field.
We prove it using rigidified Picard norms on curves.
It is used, through Corollary~\ref{cor:paper-congruence-splitting}, in Appendix~\ref{app:second-tc}, and Remark~\ref{rem:paper-local-completion} records a further application; Section~\ref{sec:paper-Q-local-global} uses the reciprocity on curves and a global analogue of the Picard argument.
Throughout this subsection, $F$ is a field and $D$ is a central division algebra over $F$.

\begin{definition}\label{def:paper-transfer}
For a central division algebra $D$ over $F$, we define the splitting transfer subgroup $I_F(D)$, the splitting transfer quotient $Q_F(D)$, and the norm subgroup $N_2([D])$ by
\[
 \begin{split}
 I_F(D)&=\sum_{D_E\,\mathrm{split}}\theta_E K_2(E),\\
 Q_F(D)&=K_2(D)/I_F(D),\\
 N_2([D])&=\sum_{D_E\,\mathrm{split}}N_{E/F}K_2(E)\subseteq K_2(F).
 \end{split}
\]
For any $\alpha\in\Br(F)$, set $N_2(\alpha)=\sum_{\alpha_E=0}N_{E/F}K_2(E)$.
In all these sums, $E/F$ ranges over finite extensions, including inseparable ones.
\end{definition}

For a field extension $L/F$ we write $\res_{L/F}\colon K(D)\to K(D_L)$ for extension of scalars and, if $L/F$ is finite, $\cor_{L/F}=N_{D_L/D}\colon K(D_L)\to K(D)$.
If $E/L$ is a finite splitting field of $D_L$, we write $\theta_{E/L}\colon K(E)\to K(D_L)$ for the splitting transfer of $D_L$, normalized as in Section~\ref{sec:paper-local-norm}; thus $\theta_E=\theta_{E/F}$ and $\cor_{L/F}\circ\theta_{E/L}=\theta_{E/F}$.
With these conventions, $I_F(A)$ and $Q_F(A)$ are defined for every central simple algebra $A$ over $F$.

The transfer formula $\Nrd_2\theta_E=N_{E/F}$ gives $N_2([D])=\Nrd_2(I_F(D))$.
With $C_2(D/F)=\coker(\Nrd_2)$, the snake lemma for the reduced norm and its restriction $I_F(D)\to N_2([D])$ gives the exact sequence
\[
 0\longrightarrow\SK_2(D)\cap I_F(D)\longrightarrow\SK_2(D)
 \longrightarrow Q_F(D)\longrightarrow K_2(F)/N_2([D])
 \longrightarrow C_2(D/F)\longrightarrow0.
\]


\begin{lemma}\label{lem:paper-Q-functorial}
Let $F'/F$ be a field extension.
Then restriction maps $I_F(D)$ into $I_{F'}(D_{F'})$, where the splitting transfers for the central simple algebra $D_{F'}$ are normalized by Morita equivalence as in Section~\ref{sec:paper-local-norm}; hence it induces a map $Q_F(D)\to Q_{F'}(D_{F'})$.
Moreover,
\[
 \deg D\cdot Q_F(D)=0.
\]
\end{lemma}
\begin{proof}
Let $E/F$ be a finite splitting field, and write
\[
 E\otimes_F F'=\prod_i R_i,
\]
where $\kappa_i$ and $\mathfrak m_i$ are the residue field and the radical of the Artinian local ring $R_i$.
Filtering restriction of scalars after base change by the powers $\mathfrak m_i^j$ on each factor, we obtain as $j$th graded piece the composite of restriction to $\kappa_i$ and splitting transfer from $\kappa_i$ to $F'$, with multiplicity $\dim_{\kappa_i}(\mathfrak m_i^j/\mathfrak m_i^{j+1})$.
The simple module of the split algebra remains simple after base change to $\kappa_i$, so the graded pieces respect the Morita normalization.
Additivity \cite{QuillenHigher} therefore gives
\[
 \res_{F'/F}\theta_E
  =\sum_i \operatorname{length}_{R_i}(R_i)\,
          \theta_{\kappa_i/F'}\,\res_{\kappa_i/E},
\]
which proves the first assertion; the lengths may exceed one when $E/F$ is inseparable.

For the second assertion, let $L\subset D$ be a separable maximal subfield, so that $[L:F]=\deg D$ and $L$ splits $D$.
For $x\in K_2(D)$ we have
\[
 \deg D\cdot x=N_{D_L/D}\bigl(\res_{L/F}x\bigr)\in\theta_LK_2(L)\subseteq I_F(D),
\]
since $K(D_L)$ is identified with $K(L)$ and $\theta_L=N_{D_L/D}$ under this identification \cite{KahnLevine2010}.
\end{proof}

The lemma on transfers of mixed products is used in part~(1) of the reciprocity law on curves.
Part~(2), for functions with residue one along the boundary, is used in Proposition~\ref{prop:paper-picard-eval}.

\begin{lemma}\label{lem:paper-mixed-products}
If $F$ is infinite, then for every finite extension $L/F$, the transfer to $F$ maps the mixed products $L^\times\cdot K_1(D_L)\subset K_2(D_L)$ into $I_F(D)$.
\end{lemma}

\begin{proof}
Let $D'$ be the division algebra in the Brauer class of $D_L$.
By Morita equivalence, we represent an element of $K_1(D_L)$ by $b\in D'^\times$.
Since $L$ is infinite, we can choose $c$ such that both $c$ and $cb$ lie in the nonempty Zariski open set of units generating separable maximal subfields.
Each factor of $b=c^{-1}\cdot(cb)$ then has the form $\theta_{E/L}(u)$ in $K_1(D_L)$ for a separable maximal subfield $E\subset D'$ and some $u\in E^\times$.
For $a\in L^\times$, the projection formula and transitivity of transfer give
\[
 \cor_{L/F}\bigl(\{a\}\cdot\theta_{E/L}(u)\bigr)
 =\cor_{L/F}\theta_{E/L}(\{a|_E,u\})=\theta_E(\{a|_E,u\}),
\]
which lies in $I_F(D)$ since $E$ splits $D$.
By bilinearity, the transfer of $\{a\}\cdot b$ to $F$ lies in $I_F(D)$.
\end{proof}

\begin{lemma}\label{lem:paper-curve-reciprocity}
Let $Y/F$ be a proper geometrically integral curve, let $V\subset Y$ be a smooth open subset, and let $\alpha\in K_2(V;D)$, where $K(V;D)$ is the $K$-theory of the Azumaya algebra $\mathcal O_V\otimes_FD$ over $V$.
Let $f$ be a rational function on the normalization of $Y$ that is a unit at every branch of the boundary $Y\setminus V$, and put
\[
 \Sigma(f)=\sum_{z\in V\ \mathrm{closed}}v_z(f)\,\cor_{\kappa(z)/F}\bigl(\alpha(z)\bigr)\in K_2(D).
\]
\begin{enumerate}
\item If $F$ is infinite, then $\Sigma(f)\in I_F(D)$.
\item If $f$ has residue $1$ at every branch of the boundary, then $\Sigma(f)=0$.
\item Suppose that $F$ is infinite, and let $P,P'\in V(F)$ satisfy $P-P'=\operatorname{div}(f)+Z$ on the normalization, where $f$ is as above and $Z$ is supported on closed points of $V$ whose residue fields split $D$.
Then $[\alpha(P)]=[\alpha(P')]$ in $Q_F(D)$.
\end{enumerate}
\end{lemma}
\begin{proof}
On the normalization of $Y$, which is unchanged over $V$, we form $f\cdot\alpha$ at the generic point.
The composite of the localization boundary and proper pushforward is zero \cite{QuillenHigher}.
At a closed point $z$ of $V$, the contribution is $v_z(f)\cor_{\kappa(z)/F}(\alpha(z))$.
At the boundary, the contributions are mixed products of the residue $\overline f$ with the $K_1$ boundary of $\alpha$; they lie in $I_F(D)$ by Lemma~\ref{lem:paper-mixed-products} if $F$ is infinite, which gives~(1), and they vanish if $\overline f=1$ at every branch, which gives~(2).
Assertion~(3) follows from~(1), since the contributions of split closed points lie in $I_F(D)$.
\end{proof}

\begin{proposition}\label{prop:paper-picard-eval}
Let $F$ be complete for a nontrivial rank-$1$ nonarchimedean valuation, and let $E/F$ be a finite separable splitting field of $D$.
For a smooth affine $F$-scheme $X$ and $\alpha\in K_2(X;D)$, the map
\[
 X(F)\longrightarrow K_2(D)/\theta_E K_2(E),\qquad P\longmapsto[P^*\alpha],
\]
is locally constant.
\end{proposition}

We change one coordinate at a time and compare the values on a curve using the rigidified Picard norm and Lemma~\ref{lem:paper-curve-reciprocity}(2).
The rigidification along the boundary makes the boundary contributions vanish, which gives constancy modulo the single subgroup $\theta_EK_2(E)$ rather than modulo $I_F(D)$.

\begin{proof}
We first set up a family of curves with rigidified Picard spaces.
In the projective closure of an affine embedding of $X$, let $H_\infty$ be the hyperplane at infinity.
We take a projective flat family $Y\to B$ of geometrically integral curves, parametrizing the complete-intersection curves constructed in the last paragraph, meeting $X$ in every fiber.
The boundary $M=Y\cap H_\infty$ is a finite flat relative Cartier divisor over $B$.
The forgetful map of the relative Picard space $\mathcal P_M$ with trivializations along $M$ has kernel
\[
 \operatorname{Res}_{M/B}\G_m/\G_m,
\]
which is smooth because it is the open locus of units in a relative projective bundle.
The Picard space is smooth since $H^2(\mathcal O)$ vanishes for curves.
On each proper geometrically integral fiber, regular functions come from the base, so a rigidified line bundle has trivial automorphism group, since $M$ is a nonempty finite flat boundary; hence $F$-points of $\mathcal P_M$ are represented by rigidified line bundles on the fibers.

Next we solve the norm equation near the diagonal.
The Picard norm along $E$ is smooth because it becomes addition after finite \'etale splitting.
For two points $P$ and $P'$ on the smooth locus, the line bundle $\mathcal O(P-P')$ has a canonical trivialization along $M$.
Starting from the trivial solution along the diagonal, we obtain nearby solutions by the \'etale inverse function theorem, which follows here from Hensel's lemma for a complete rank-$1$ valued field.
Let $\mathcal M$ be such a solution: a line bundle on $Y_E$, rigidified along $M$, whose norm is $\mathcal O(P-P')$ with its canonical rigidification.
Since line bundles over semilocal rings are free and units lift from the finite boundary, we can choose a rational section of $\mathcal M$ that takes the value $1$ along the boundary under the rigidification.
If $Z$ is the divisor of the section, then on the normalization
\[
 P-P'-\operatorname{Nm}(Z)=\operatorname{div}(f),\qquad f|_M=1.
\]
We also choose the section to be invertible at the points where the affine curve is not smooth over $F$, so that the support of $\operatorname{Nm}(Z)$ lies in the smooth open subset.
At the branches over the non-smooth points of the affine curve, $f$ is a unit, and $\alpha$ extends along the pullback to the normalization, so its $K_1$ boundary vanishes; at the branches at infinity, the rigidification makes the residue of $f$ equal to $1$.
Thus the boundary terms vanish by the localization calculation in the proof of Lemma~\ref{lem:paper-curve-reciprocity}, and the remaining terms are transfers through $E$.
Indeed, for a closed point $w$ in the support of $Z$ with image $z$, the pushforward of $[w]$ contributes $[\kappa(w):\kappa(z)]\cdot\cor_{\kappa(z)/F}\alpha(z)=\cor_{\kappa(w)/F}\alpha(w)$, which factors through $\theta_E$ since $E\subset\kappa(w)$.

Finally, we apply this calculation near an arbitrary smooth point.
The connected component of $X$ containing that point is smooth with an $F$-point, hence geometrically integral, and we replace $X$ by it; if it has dimension zero, local constancy is immediate.
By Bertini's theorem we choose geometrically integral complete-intersection curves through that point tangent to the coordinate axes; sections of sufficiently large degree remain free away from the point when the point and the tangent direction are prescribed.
For two nearby points differing only in one \'etale coordinate $t_i=r_i/s_0$, we correct each defining section by multiples of $s_0$ and $r_i$ so that it vanishes at both points.
Since the central curve is tangent to the $i$th axis, the correction is small, so flatness, geometric integrality, and the open norm condition persist.
In a product coordinate neighborhood, we move from one point to another by changing one coordinate at a time and apply the preceding calculation at each step.
\end{proof}

To apply local constancy to Steinberg words, we use the universal $K_2$-class on the smooth locus of the identity fiber of the word map defined below.
Denote the standard generators of the stable Steinberg group by $x_{ij}(a)$.
For a stable Steinberg word over $D$, we use an independent affine variable for each occurrence of a coefficient and take the matrix product in $\operatorname{SL}_n(D)$, the algebraic group of elements of reduced norm one in $M_n(D)$, for sufficiently large $n$; we call the resulting morphism the word map.

\begin{lemma}\label{lem:paper-steinberg-fiber}
Append to a stable Steinberg word over $D$ the letters
\[
 x_{ij}(t)\qquad\text{and}\qquad x_{ij}(u)\,x_{ji}(t)\,x_{ij}(-u),
\]
for all $i\ne j$ and all $u$ in a basis of $D$ containing $1$, with an independent $D$-valued variable $t$ for each of these letters.
Then the identity fiber of the word map is smooth at every point where the added variables are zero, and on the smooth open locus of the identity fiber the universal Steinberg word defines a $K_2$-class.
Moreover, the word map has an \'etale local section near the identity, sending the identity to the point at which the original variable coefficients and the added variables $t$ vanish.
\end{lemma}
\begin{proof}
At $t=0$ the word is unchanged.
The differentials of the added letters generate the off-diagonal entries and the elements $ute_{ii}-tue_{jj}$.
The diagonal differences for $u=1$, together with the linear span of the commutators $ut-tu$, give the entire subspace of reduced trace zero, since by nondegeneracy of the reduced trace pairing the linear span of the commutators in $D$ is $\ker\operatorname{Trd}$, also when the characteristic divides $\deg D$.
Thus the differential of the matrix product is surjective where the added variables vanish, and the identity fiber is smooth there.
On its smooth open locus, the universal Steinberg word has matrix product one and defines a $K_2$-class.
The same surjectivity gives an \'etale local section of the word map near the identity through the point where the original variable coefficients and the added variables $t$ vanish.
There the fixed letters $x_{ij}(u)$ and $x_{ij}(-u)$ cancel in the Steinberg group, so the word is trivial.
\end{proof}

\begin{corollary}\label{cor:paper-congruence-splitting}
Let $F$ be complete for a nontrivial rank-$1$ nonarchimedean valuation, and let $E/F$ be a finite separable splitting field of $D$.
For $n$ sufficiently large, let $\widetilde G\to E_n(D)$ be the pullback of the stable universal central extension by $K_2(D)$ to the elementary subgroup $E_n(D)$ of $\operatorname{SL}_n(D)$, and push it out to $K_2(D)/\theta_E K_2(E)$.
For $\gamma>0$ put $J_\gamma=\{x\in D:v_D(x)>\gamma\}$, where $v_D$ is the valuation of $D$.
Then this central extension splits over $E_n(D)\cap(1+\operatorname{Mat}_n(J_\gamma))$ for all sufficiently large $\gamma$.
In the discrete case, this gives a splitting over $E_n(D)\cap(1+\pi_F^m\operatorname{Mat}_n(\mathcal O_D))$ for all sufficiently large $m$, where $\pi_F$ is a uniformizer of $F$ and $\mathcal O_D$ is the valuation ring of $D$.
\end{corollary}
\begin{proof}
We choose an affine \'etale neighborhood $\rho\colon V\to\operatorname{SL}_n(D)$ of the identity on which the \'etale local section of Lemma~\ref{lem:paper-steinberg-fiber} (applied to the empty word) is defined and sends the identity to the trivial word, and we lift points near the identity by Hensel's lemma.
The multiplicative error of this section is a universal $K_2$-class on the closed subscheme of $V^3$ where $\rho(c)=\rho(a)\rho(b)$, which is smooth and affine near the identity triple, and we apply Proposition~\ref{prop:paper-picard-eval} to it.
The error vanishes at the identity triple, and hence, by local constancy, on a neighborhood of it.
Choosing $\gamma$ so that the product of the subgroup $E_n(D)\cap(1+\operatorname{Mat}_n(J_\gamma))$ with itself lies in this neighborhood, we obtain the splitting over this subgroup, and hence over its subgroups for larger $\gamma$.
\end{proof}

\begin{remark}\label{rem:paper-local-completion}
We have also proved the following extension of Theorem~\ref{thm:paper-local-norm} in degree two; it is not used in this paper, and we omit the proof.
Let $H$ be an algebraic extension of a nonarchimedean local field, and let $\widehat H$ be its completion.
Then, for every central division algebra $D$ over $\widehat H$, the reduced norm $\Nrd_2\colon K_2(D)\to K_2(\widehat H)$ is an isomorphism.
The proof descends $D$ to a finite subextension and applies Proposition~\ref{prop:paper-picard-eval} over $\widehat H$.
\end{remark}
